\chapter*{Coda}

% Phase 4 scaffold. Brief return to the road — the Toronto classroom scene
% with the model from ch 13 now visible. Echoes ch 01's opening (the road
% is long long).
%
% Length target: 3-5 manuscript pages.
%
% Per Zimmer's arc proposal: "Resolution echoes opening: back to the road,
% the students, the songwriter. Takeaway — grammaticality is not what you
% thought it was, and that is a relief, not a problem."
%
% Per Morris's arc proposal: "Last shot returns to the classroom, asterisk
% erased."
%
% Suggested structure (NO sections; flow as continuous prose):
%   - Open in the classroom again. Same Turkish student. Same "the road
%     is long long."
%   - This time the reader (and you) can see the cluster — the form sits
%     in a basin English speakers don't ordinarily inhabit but which is
%     close enough to neighbours (intensifying reduplication, contact
%     features, register effects) that it lands as creative rather than
%     wrong.
%   - The asterisk is now a measurement, not a judgement.
%   - The book ends not on definition but on the act of perception.
%
% Use \chapter* (unnumbered) so the coda isn't numbered as ch 15 in the
% TOC — it functions as an envoi rather than a chapter.
